%------------------------------------------------------------------------------%
\begin{question}
  Descreva o domínio da função
  \(
    f(x, y) = \frac{\sqrt{x^2-4\,}}{\;\ln(y-1)\;}
  \)
  \begin{enumerate}
    \item Qual a condição que os pontos do domínio deve satisfazer?
    \item Que região é essa? Represente-a graficamente.
    \item Qual o interior da região?
    \item Qual a fronteira da região?
    \item A região é aberta ou fechada?
    \item A região é limitada?
  \end{enumerate}
\end{question}

%------------------------------------------------------------------------------%
\begin{resolution}{Solução}
  Condições para avaliar a função
  \(f(x, y) = \frac{\sqrt{x^2-4\,}}{\;\ln(y-1)\;}\)

  \pause
  \begin{enumerate}[<+->]
    \item \(x^2-4 \geq 0\)
    \item \(\ln(y-1) \neq 0\)
    \item \(y-1 > 0\)
  \end{enumerate}
\end{resolution}

%------------------------------------------------------------------------------%
\begin{resolution}{Solução}
\begin{columns}[t]
\column{0.3\linewidth}
  \onslide<+->{Condição 1}
  \begin{align*}
    \onslide<+->{x^2-4     & \geq 0                       \\}
    \onslide<+->{x^2       & \geq 4                       \\}
    \onslide<+->{\abs{x}   & \geq 2                       \\}
    \onslide<+->{x \leq -2 & \quad\text{ou}\quad x \geq 2}
  \end{align*}
\column{0.3\linewidth}
  \onslide<+->{Condição 2}
  \begin{align*}
    \onslide<+->{\ln(y-1) & \neq 0 \\}
    \onslide<+->{y-1      & \neq 1 \\}
    \onslide<+->{y        & \neq 2}
  \end{align*}
\column{0.3\linewidth}
  \onslide<+->{Condição 3}
  \begin{align*}
    \onslide<+->{y-1 & > 0 \\}
    \onslide<+->{y   & > 1}
  \end{align*}
\end{columns}
\end{resolution}

%------------------------------------------------------------------------------%
\begin{resolution}{Solução}
  Domínio
  \[
    D_f = \left\{ (x, y)\in\R^2 \st
      \abs{x} \geq 2 \;\text{ e }\;
      y > 1          \;\text{ e }\;
      y \neq 2
    \right\}
  \]
\end{resolution}

%------------------------------------------------------------------------------%
\begin{resolution}{Representação Gráfica}
  \centering
  \includegraphics{ex_dominio_sqrt_x2_4_ln_y_1}
\end{resolution}

%------------------------------------------------------------------------------%
\begin{resolution}{Solução}
  \begin{enumerate}
    \item \onslide<+->{Interior da região}
    \onslide<+->{
    \[
      \text{Interior } = \left\{ (x, y)\in\R^2 \st
          \abs{x} > 2 \;\text{ e }\;
          y > 1       \;\text{ e }\;
          y \neq 2
        \right\}
    \]}
    \vspace{-\baselineskip}
    \item \onslide<+->{Fronteira da região}
    \onslide<+->{
    \begin{align*}
      \text{Fronteira } = \Big\{ (x, y)\in\R^2 \st
          & y = 1 \text{ e } \abs{x} \geq 2, \text{ ou }\\[-2mm]
          & y = 2 \text{ e } \abs{x} \geq 2, \text{ ou }
          x = \pm 2 \text{ e } y   \geq 1
        \Big\}
    \end{align*}}
    \vspace{-\baselineskip}
    \item \onslide<+->A região é aberta ou fechada? \\
    \onslide<+->{Nem aberta nem fechada}
    \item \onslide<+->{A região é limitada?} \\
    \onslide<+->{Ilimitada,}
    \onslide<+->{pois dado qualquer disco sempre teremos um ponto fora dele}
  \end{enumerate}
\end{resolution}
%------------------------------------------------------------------------------%